\section{Análisis GAP}
\label{sec:GapAnalysis}

El análisis se ha estructurado siguiendo el enfoque propuesto por el \textit{International Institute of Business Analysis}, basado en comparar el estado actual con el estado objetivo, identificar las brechas relevantes y documentar recomendaciones o respuestas de actuación \citep{iiba-gap-analysis}. En este caso, dicho enfoque se traslada a una matriz práctica formada por cinco elementos: dimensión funcional, situación actual, estado objetivo, brecha identificada y respuesta del proyecto.

El análisis se centra en las \textbf{capacidades que debe alcanzar el sistema}: gestión de usuarios, clientes profesionales, catálogo técnico, visitas comerciales, rutas, pedidos, repartos, cobros, comunicaciones, trazabilidad y acceso diferenciado por roles. A partir de esa comparación, se identifica para cada dimensión la respuesta funcional prevista.

La Tabla~\ref{tab:gap-analysis} organiza el análisis por dimensiones funcionales. Para cada una de ellas se recoge la situación actual, el estado objetivo perseguido, la brecha identificada y la respuesta planteada por el proyecto.

{\scriptsize
\setlength{\parskip}{0pt}
\setlength{\tabcolsep}{2.5pt}
\renewcommand{\arraystretch}{1.04}
\tablecaption{Matriz GAP entre situación actual, estado objetivo y respuesta del sistema\label{tab:gap-analysis}}
\tablefirsthead{%
\hline
\rowcolor{black!8}
\textbf{Dimensión} & \textbf{Situación actual} & \textbf{Estado objetivo} & \textbf{Brecha identificada} & \textbf{Respuesta del proyecto} \\
\hline}
\tablehead{%
\hline
\rowcolor{black!8}
\textbf{Dimensión} & \textbf{Situación actual} & \textbf{Estado objetivo} & \textbf{Brecha identificada} & \textbf{Respuesta del proyecto} \\
\hline}

\begin{supertabular}{|>{\raggedright\arraybackslash}p{0.16\linewidth}|>{\raggedright\arraybackslash}p{0.19\linewidth}|>{\raggedright\arraybackslash}p{0.20\linewidth}|>{\raggedright\arraybackslash}p{0.20\linewidth}|>{\raggedright\arraybackslash}p{0.16\linewidth}|}

\textbf{Usuarios, roles y seguridad} & El acceso a la información depende de controles manuales, conocimiento interno y herramientas no integradas. No existe un circuito único para altas, activaciones, bloqueos o permisos. & Disponer de un sistema con \textbf{autenticación}, revisión de altas, estados de cuenta, roles diferenciados y trazabilidad básica de accesos. & Falta un mecanismo centralizado que controle quién accede, qué puede hacer y cómo se auditan las acciones relevantes. & M1 agrupa el acceso, el alta y la seguridad de cuentas; se apoya en requisitos funcionales, no funcionales y reglas de negocio de control de acceso. \\
\hline

\textbf{Clientes profesionales y salones} & La información de peluquerías, responsables, datos técnicos, condiciones comerciales y relación con comerciales se encuentra dispersa entre conversaciones, documentos y conocimiento operativo. & Mantener una ficha única del cliente profesional, con datos comerciales, técnicos, de contacto, asignaciones y estado administrativo. & No hay una visión consolidada del cliente que conecte la actividad comercial, la información técnica del salón y la gestión administrativa. & M2 centraliza la gestión de clientes profesionales y salones, enlazando requisitos funcionales y de información sobre clientes, contactos y datos técnicos. \\
\hline

\textbf{Catálogo técnico y productos} & El catálogo y la información técnica se consultan mediante canales externos, catálogos de marca o documentación no conectada con la operativa diaria del distribuidor. & Gestionar un catálogo propio con productos, categorías, imágenes, referencias técnicas y cartas de color vinculadas al uso comercial. & Las soluciones existentes aportan consulta de producto, pero no conectan el catálogo con pedidos, promociones, roles y gestión local del distribuidor. & M3 agrupa catálogo, productos, familias, cartas de color y recursos gráficos necesarios para la venta y la consulta técnica. \\
\hline

\textbf{Visitas comerciales y rutas} & La planificación de visitas y desplazamientos se realiza de forma operativa, apoyada en criterio del comercial, acuerdos puntuales y herramientas de mapas externas. & Registrar visitas, organizar rutas, consultar ubicaciones, asignar comerciales y facilitar una planificación diaria coherente con clientes y pedidos. & La capa geográfica y la actividad comercial no están conectadas con la información de clientes, pedidos y seguimiento del trabajo realizado. & M2 concentra la gestión comercial, las visitas, la planificación de rutas, la geolocalización y el apoyo a la actividad presencial del comercial. \\
\hline

\textbf{Pedidos, repartos y cobros} & Los pedidos, entregas y cobros forman parte de un flujo operativo que puede quedar repartido entre mensajes, documentos, anotaciones y sistemas administrativos externos. & Integrar el ciclo desde la solicitud del pedido hasta la preparación, reparto, confirmación de entrega y seguimiento económico. & No existe continuidad digital entre venta, logística y control económico, lo que dificulta saber el estado real de cada operación. & M4 cubre pedidos, preparación, reparto, validación mediante \textit{QR}, estados y control operativo de cobros. \\
\hline

\textbf{Comunicaciones y promociones} & Las comunicaciones comerciales, avisos y promociones pueden depender de canales externos o intercambios individuales, sin trazabilidad uniforme dentro del sistema. & Disponer de avisos, notificaciones, promociones y mensajes operativos asociados a usuarios, clientes, pedidos o campañas concretas. & La comunicación no queda unida al dato de negocio ni al historial de interacción con el cliente profesional. & M6 agrupa promociones, notificaciones, comunicaciones y formaciones, apoyándose en M7 para la configuración transversal cuando sea necesario. \\
\hline

\textbf{Uso móvil y experiencia unificada} & Los distintos procesos pueden apoyarse en herramientas independientes que no ofrecen una experiencia común para administrador, comercial y cliente profesional. & Ofrecer una aplicación web responsive, con base \textit{PWA}, accesible desde escritorio y móvil para los distintos perfiles de usuario. & La movilidad existe en herramientas concretas, pero no como experiencia integrada alrededor del dominio completo del distribuidor. & Los requisitos no funcionales y la arquitectura web del sistema establecen la base de uso multiplataforma y acceso desde dispositivos móviles. \\
\hline

\textbf{Integraciones y evolución futura} & Sistemas externos como Factusol, automatizaciones específicas, asistentes inteligentes o simulación capilar no forman parte de un flujo integrado inicial. & Preparar el sistema para evolucionar hacia integraciones empresariales, automatización de procesos y funcionalidades inteligentes vinculadas al dominio capilar. & Algunas necesidades aportan valor, pero exceden el alcance principal del primer sistema y deben tratarse como evolución planificada. & M8 documenta el trabajo futuro: integración con Factusol, automatización, asistente virtual, diagnóstico con capilógrafo y simulación mediante cámara móvil. \\
\hline
\end{supertabular}
}

El resultado del análisis muestra que la brecha principal no se encuentra en una funcionalidad aislada, sino en la falta de continuidad entre tres dimensiones: la \textbf{gestión comercial interna del distribuidor}, la \textbf{experiencia digital del cliente profesional} y el \textbf{seguimiento logístico y económico} de pedidos y repartos. Las herramientas revisadas en el estado del arte pueden cubrir partes concretas de esa cadena, pero no resuelven de forma conjunta el flujo operativo completo.

Por tanto, el sistema propuesto se justifica como una \textbf{solución específica} y no como una simple sustitución de una hoja de cálculo, un catálogo digital, un \textit{CRM}, un \textit{ERP} o una herramienta de rutas. El análisis orienta la solución hacia la centralización del dominio del distribuidor sobre una misma base de datos, con \textbf{roles diferenciados}, \textbf{trazabilidad operativa}, soporte \textit{PWA} \citep{webdev-pwa} y una estructura modular preparada para integrar en el futuro herramientas externas como Factusol, flujos con \textit{n8n} o servicios especializados de automatización \citep{sdelsol-factusol,n8n-webhook,n8n-http-request}.
